\chapter*{Danh mục từ viết tắt và thuật ngữ}
\addcontentsline{toc}{chapter}{Danh mục từ viết tắt và thuật ngữ}

{\renewcommand{\arraystretch}{1.25}
\noindent\begin{tabularx}{\textwidth}{@{}>{\bfseries}p{2.2cm}>{\raggedright\arraybackslash}X>{\raggedright\arraybackslash}X@{}}
\toprule
Viết tắt & \textbf{Tiếng Anh} & \textbf{Tiếng Việt}\\
\midrule
ARI & Adjusted Rand Index & Chỉ số Rand hiệu chỉnh\\
Bi-GRU & Bidirectional Gated Recurrent Unit & Mạng hồi quy có cổng hai chiều\\
CASAS & Center for Advanced Studies in Adaptive Systems & Dự án nhà thông minh của Đại học Bang Washington\\
CPD & Change Point Detection & Phát hiện điểm thay đổi\\
DBSCAN & Density-Based Spatial Clustering of Applications with Noise & Phân cụm theo mật độ có nhiễu\\
F1 & F1-score & Trung bình điều hòa của độ chính xác và độ phủ\\
HAR & Human Activity Recognition & Nhận dạng hoạt động của con người\\
HDBSCAN & Hierarchical DBSCAN & DBSCAN phân cấp\\
LLM & Large Language Model & Mô hình ngôn ngữ lớn\\
P, R & Precision, Recall & Độ chính xác, độ phủ\\
PIR & Passive Infrared & Cảm biến hồng ngoại thụ động (chuyển động)\\
RF & Random Forest & Rừng ngẫu nhiên\\
RQ & Research Question & Câu hỏi nghiên cứu\\
\bottomrule
\end{tabularx}\par}

\vspace{0.8cm}
\noindent{\large\bfseries Một số thuật ngữ dùng trong khóa luận}
\begin{itemize}
  \item \textbf{Sự kiện cảm biến}: một dòng nhật ký ``thời điểm $t$, cảm biến $s$, giá trị $v$''.
  \item \textbf{Lần thực hiện} (episode): một đoạn thời gian $[t^{s}, t^{e}]$ trong đó một người thực hiện
        một hoạt động.
  \item \textbf{Thói quen}: tập các lần thực hiện cùng một hoạt động, bởi cùng một người, bắt đầu quanh
        cùng một giờ, lặp lại ở ít nhất một nửa số ngày (Định nghĩa~\ref{def:habit}).
  \item \textbf{Thói quen đáp án}: thói quen thu được khi áp định nghĩa lên nhãn thật; chỉ dùng để đánh giá.
  \item \textbf{Vùng}: nhóm cảm biến gần nhau, rút ra từ dữ liệu, đóng vai trò ``phòng''.
  \item \textbf{Vệt} (track): chuỗi sự kiện được cho là do cùng một người gây ra.
  \item \textbf{Nhãn thật (Oracle)}: khai phá thói quen trực tiếp trên nhãn thật; là giới hạn trên của
        mọi phương pháp.
\end{itemize}
